\documentclass[conference]{IEEEtran}
\IEEEoverridecommandlockouts
\usepackage{cite}
\usepackage{amsmath,amssymb}
\usepackage{graphicx}
\usepackage{booktabs}
\usepackage{url}
\usepackage{textcomp}
\usepackage{xcolor}
\newcommand{\nQ}{100}
\newcommand{\nPar}{991}
\newcommand{\mainModel}{gpt-4o-mini}
\newcommand{\fastModel}{gpt-4o-mini}
\newcommand{\retrKind}{tfidf}
\newcommand{\emScb}{0.0}
\newcommand{\fOnecb}{0.0}
\newcommand{\fLocb}{0.0}
\newcommand{\fHicb}{0.0}
\newcommand{\ctxTokcb}{0}
\newcommand{\totTokcb}{416}
\newcommand{\callscb}{2.00}
\newcommand{\recallcb}{0.0}
\newcommand{\latcb}{1.54}
\newcommand{\iterscb}{0.00}
\newcommand{\violcb}{2.000}
\newcommand{\directcb}{100.0}
\newcommand{\faithcb}{--}
\newcommand{\emSnv}{25.0}
\newcommand{\fOnenv}{34.1}
\newcommand{\fLonv}{25.4}
\newcommand{\fHinv}{43.0}
\newcommand{\ctxToknv}{466}
\newcommand{\totToknv}{956}
\newcommand{\callsnv}{1.40}
\newcommand{\recallnv}{75.5}
\newcommand{\latnv}{1.31}
\newcommand{\itersnv}{0.00}
\newcommand{\violnv}{0.800}
\newcommand{\directnv}{0.0}
\newcommand{\faithnv}{--}
\newcommand{\emSit}{58.0}
\newcommand{\fOneit}{74.2}
\newcommand{\fLoit}{66.5}
\newcommand{\fHiit}{81.6}
\newcommand{\ctxTokit}{575}
\newcommand{\totTokit}{2660}
\newcommand{\callsit}{3.55}
\newcommand{\recallit}{91.0}
\newcommand{\latit}{3.58}
\newcommand{\itersit}{0.88}
\newcommand{\violit}{0.000}
\newcommand{\directit}{0.0}
\newcommand{\faithit}{--}
\newcommand{\emSar}{54.0}
\newcommand{\fOnear}{68.4}
\newcommand{\fLoar}{61.1}
\newcommand{\fHiar}{76.7}
\newcommand{\ctxTokar}{661}
\newcommand{\totTokar}{3213}
\newcommand{\callsar}{4.30}
\newcommand{\recallar}{83.0}
\newcommand{\latar}{4.63}
\newcommand{\itersar}{0.79}
\newcommand{\violar}{0.000}
\newcommand{\directar}{8.0}
\newcommand{\faithar}{--}
\newcommand{\emScr}{55.0}
\newcommand{\fOnecr}{68.9}
\newcommand{\fLocr}{60.8}
\newcommand{\fHicr}{76.7}
\newcommand{\ctxTokcr}{544}
\newcommand{\totTokcr}{2752}
\newcommand{\callscr}{4.27}
\newcommand{\recallcr}{82.5}
\newcommand{\latcr}{4.52}
\newcommand{\iterscr}{0.79}
\newcommand{\violcr}{0.000}
\newcommand{\directcr}{8.0}
\newcommand{\faithcr}{--}
\newcommand{\emScrfb}{55.0}
\newcommand{\fOnecrfb}{70.2}
\newcommand{\fLocrfb}{62.2}
\newcommand{\fHicrfb}{77.9}
\newcommand{\ctxTokcrfb}{748}
\newcommand{\totTokcrfb}{3345}
\newcommand{\callscrfb}{4.10}
\newcommand{\recallcrfb}{87.5}
\newcommand{\latcrfb}{4.34}
\newcommand{\iterscrfb}{0.69}
\newcommand{\violcrfb}{0.000}
\newcommand{\directcrfb}{8.0}
\newcommand{\faithcrfb}{--}
\newcommand{\emScrnc}{45.0}
\newcommand{\fOnecrnc}{60.0}
\newcommand{\fLocrnc}{51.5}
\newcommand{\fHicrnc}{68.9}
\newcommand{\ctxTokcrnc}{530}
\newcommand{\totTokcrnc}{886}
\newcommand{\callscrnc}{2.00}
\newcommand{\recallcrnc}{76.5}
\newcommand{\latcrnc}{2.21}
\newcommand{\iterscrnc}{0.00}
\newcommand{\violcrnc}{0.000}
\newcommand{\directcrnc}{8.0}
\newcommand{\faithcrnc}{--}
\newcommand{\dFar}{+0.5}
\newcommand{\dTokPctar}{-17.7}
\newcommand{\dTotPctar}{-14.4}
\newcommand{\dCallsPctar}{-0.7}
\newcommand{\dFloar}{-3.4}
\newcommand{\dFhiar}{+4.7}
\newcommand{\dFnv}{+34.8}
\newcommand{\dTokPctnv}{+16.9}
\newcommand{\dTotPctnv}{+187.7}
\newcommand{\dCallsPctnv}{+205.0}
\newcommand{\dFlonv}{+25.3}
\newcommand{\dFhinv}{+44.5}
\newcommand{\dFit}{-5.3}
\newcommand{\dTokPctit}{-5.4}
\newcommand{\dTotPctit}{+3.4}
\newcommand{\dCallsPctit}{+20.3}
\newcommand{\dFloit}{-11.0}
\newcommand{\dFhiit}{-0.3}
\newcommand{\swLowB}{150}
\newcommand{\swMidB}{600}
\newcommand{\swcrLow}{37.7}
\newcommand{\swcrMid}{60.3}
\newcommand{\swrawLow}{39.7}
\newcommand{\swrawMid}{57.5}
\newcommand{\swWins}{3}
\newcommand{\swTotal}{5}
\newcommand{\swMeanGain}{+0.6}
   % auto-generated by scripts/make_report.py from results/*.csv
\newcommand{\figorpending}[2]{\IfFileExists{#1}{\includegraphics[width=#2]{#1}}{\fbox{figure pending: run make\_report.py}}}

\begin{document}

\title{ContextRouter: Adaptive Retrieval-Augmented Generation with Agent-Level Context-Budget Engineering and Typed Agent Contracts}

\author{\IEEEauthorblockN{Moulidharan S}
\IEEEauthorblockA{Department of Computing -- Decision and Computing Sciences\\
Coimbatore Institute of Technology, Coimbatore, India}}

\maketitle

\begin{abstract}
Adaptive retrieval-augmented generation (RAG) systems decide \emph{whether} and \emph{how often} to retrieve, but
they still fill the language model's context window with whole retrieved passages. We present ContextRouter, a
multi-agent RAG workflow in which a dedicated Context Engineer agent decides \emph{what} enters the window under a
token budget chosen per question by a Strategy Planner. A Query Analyzer routes each question (direct, single-hop,
multi-hop, comparison), a Retriever builds a multi-query candidate pool, the Context Engineer de-duplicates the pool,
scores sentences (query overlap, bridge-entity links, position), packs the best ones into the budget and orders them
to avoid the ``lost in the middle'' effect, and an Answerer--Critic loop iterates with follow-up queries when the
evidence is insufficient. All inter-agent messages are typed contracts validated at run time, and every system in our
study is the same LangGraph workflow under different configuration flags, so differences are attributable to
the flags. On \nQ\ HotpotQA (distractor) questions with \mainModel\ as the answerer, ContextRouter reaches
\fOnecr\% F1 (95\% CI [\fLocr, \fHicr]) with \ctxTokcr\ context tokens on average, versus \fOnear\% F1 and \ctxTokar\
tokens for the same adaptive pipeline without context engineering. We also report ablations, paired bootstrap
confidence intervals, and an accuracy-versus-budget sweep comparing context engineering with hard truncation. Code,
data subset and evaluation scripts are released for reproduction.
\end{abstract}

\begin{IEEEkeywords}
agentic AI, retrieval-augmented generation, context engineering, multi-hop question answering, multi-agent systems
\end{IEEEkeywords}

\section{Introduction}
Retrieval-augmented generation (RAG) grounds large language models (LLMs) in external documents \cite{lewis2020rag}.
Recent work makes RAG \emph{adaptive}: the system routes a question to no retrieval, single-step retrieval or
multi-step retrieval according to its complexity \cite{jeong2024adaptive}, or lets the model critique its own
retrievals \cite{asai2024selfrag,yan2024crag}. Yet one decision is left implicit in almost every pipeline: how the
retrieved material is assembled into the prompt. In practice whole passages are concatenated in rank order. This
spends tokens on irrelevant sentences, and long contexts are known to be used unevenly, with evidence placed in the
middle of the window being exploited less reliably \cite{liu2024lost}.

\emph{Context engineering}---deliberately curating what the model sees---is increasingly treated as a first-class design
problem for agentic systems. This paper studies it as an explicit \emph{agent responsibility} inside an adaptive RAG
workflow and asks: \textbf{can a budget-aware Context Engineer agent keep answer quality while reducing the
context an LLM must read, compared with passing whole retrieved passages?}

\textbf{Contributions.}
(1) A Context Engineer agent that packs query-relevant, bridge-linked sentences into a per-question token budget
without extra LLM calls. (2) A six-agent LangGraph workflow (Analyzer, Planner, Retriever, Context Engineer,
Answerer, Critic) coordinated through typed Pydantic contracts and a shared state, combining routing and
evaluator--optimizer patterns \cite{anthropic2024agents}. (3) A controlled evaluation in which baselines and
ablations are configurations of one workflow, with paired bootstrap confidence intervals and an accuracy-versus-budget
sweep against hard truncation. (4) An open, reproducible implementation.

\section{Related Work}
\textbf{Adaptive and corrective RAG.} Adaptive-RAG routes questions by predicted complexity \cite{jeong2024adaptive};
Self-RAG trains a model to decide when to retrieve and to critique outputs \cite{asai2024selfrag}; CRAG evaluates
retrieval quality and triggers corrective actions \cite{yan2024crag}. We keep the routing idea but add an explicit
stage controlling the \emph{content and size} of the context.

\textbf{Iterative retrieval and self-feedback.} IRCoT interleaves retrieval with chain-of-thought steps
\cite{trivedi2023ircot}; ReAct interleaves reasoning and tool use \cite{yao2023react}; Self-Refine iterates with
model feedback \cite{madaan2023selfrefine}. Our Critic follows the evaluator--optimizer pattern and issues a
follow-up retrieval query when evidence is missing.

\textbf{Context compression and long-context behavior.} LLMLingua compresses prompts with a small language model
\cite{jiang2023llmlingua}; Liu et al.\ show that models use information in the middle of long contexts less
reliably \cite{liu2024lost}. Our Context Engineer is extractive and training-free, and uses the structure of
multi-hop evidence (bridge entities) rather than token-level perplexity.

\textbf{Agent coordination.} Practical agent design distinguishes workflow patterns such as prompt chaining,
routing, parallelization, orchestrator--workers and evaluator--optimizer \cite{anthropic2024agents}. ContextRouter
composes routing and evaluator--optimizer, and makes the interfaces between agents explicit and machine-checked.

\section{Method}
\subsection{Workflow and agents}
Fig.~\ref{fig:arch} shows the workflow. Given question $q$:
\begin{enumerate}
\item \textbf{Query Analyzer} (LLM, fast model). Returns a \texttt{QueryPlan}: type
$\in\{$direct, single-hop, multi-hop, comparison$\}$ and up to three sub-questions.
\item \textbf{Strategy Planner} (policy). Maps the type to a context budget $B$ and retrieval depth $k$
(Table~\ref{tab:policy}) and forms the query set $Q=\{q\}\cup$ sub-questions. ``Direct'' questions skip retrieval.
\item \textbf{Retriever Agent.} Retrieves top-$k$ paragraphs per query ($k/2$ for sub-queries) and merges them into a
candidate pool keeping each paragraph's best score.
\item \textbf{Context Engineer} (Section~\ref{sec:ce}). Converts the pool into a context of at most $B$ tokens.
\item \textbf{Answerer} (LLM, main model). Returns an \texttt{AnswerDraft} using only the context.
\item \textbf{Critic} (LLM, fast model). Returns a \texttt{Critique} (grounded, sufficient, follow-up query). If the
evidence is insufficient and a follow-up query is given, the follow-up is retrieved, added to the pool, the context is
re-engineered under the same budget and the Answerer runs again, for at most three iterations.
\end{enumerate}

\begin{figure}[t]
\centering
\figorpending{figs/architecture.png}{\columnwidth}
\caption{ContextRouter agents, typed contracts and control flow.}
\label{fig:arch}
\end{figure}

\begin{table}[t]
\caption{Strategy-planner policy (set a priori, not tuned on evaluation data).}
\label{tab:policy}
\centering
\begin{tabular}{lcc}
\toprule
Question type & Budget $B$ (tokens) & Top-$k$ \\
\midrule
direct & 0 (no retrieval) & 0 \\
single-hop & 400 & 4 \\
multi-hop & 900 & 8 \\
comparison & 900 & 8 \\
\bottomrule
\end{tabular}
\end{table}

\subsection{Typed contracts and shared state}
Every message between agents is a Pydantic model (\texttt{QueryPlan}, \texttt{StrategyDecision},
\texttt{Chunk}, \texttt{EngineeredContext}, \texttt{AnswerDraft}, \texttt{Critique}) held in a typed LangGraph state.
An LLM reply that fails parsing or schema validation triggers one repair retry; each failure is counted as a
\emph{contract violation} and reported per system. This makes interface reliability measurable rather than assumed.

\subsection{Context Engineer}\label{sec:ce}
The agent makes no LLM calls. (i) \emph{De-duplicate}: keep each title once (best score) and drop paragraphs whose
token-set Jaccard similarity with a kept paragraph exceeds 0.8. (ii) \emph{Score sentences}. For sentence $x$ in
paragraph $p(x)$,
\begin{multline}
s(x)=w_c\,\hat r(p(x)) + w_o\,\mathrm{ov}(x,Q) + w_b\,b(x)\\
+ w_f\,\mathbb{1}[x\ \text{first}] + w_t\,t(p(x)),
\end{multline}
where $\hat r$ is the retrieval score normalised by the pool maximum; $\mathrm{ov}$ is the IDF-weighted fraction of query
terms covered by $x$; $b(x)=1$ if $x$ mentions the (parenthetical-stripped) title of another pool paragraph, which
captures bridge entities in multi-hop chains; the first sentence usually defines the entity; and $t(p)=1$ if the
paragraph title occurs in the question. We use $(w_c,w_o,w_b,w_f,w_t)=(0.5,1.0,0.7,0.3,0.4)$, fixed by hand before any
evaluation. (iii) \emph{Pack}: greedily add sentences by descending $s(x)$ while the estimated token count (1.3
tokens per word, plus paragraph titles) stays within $B$. (iv) \emph{Order}: restore sentence order within each
paragraph and arrange paragraphs so that the strongest is first and the second strongest last, mitigating
the lost-in-the-middle effect \cite{liu2024lost}.

\section{Experimental Setup}
\textbf{Data.} HotpotQA, distractor setting, validation split \cite{yang2018hotpotqa}. We sample \nQ\ questions with a fixed
random seed (42); the retrieval corpus is the union of all \nPar\ distinct paragraphs of these questions (gold plus distractors).
\textbf{Models.} Answerer: \mainModel; Analyzer and Critic: \fastModel; temperature 0, served through an
OpenAI-compatible API. Retrieval: \retrKind\ over title plus paragraph text (dense retrieval
\cite{reimers2019sbert,johnson2021faiss} is supported but not used in the main experiment).

\textbf{Systems.} All systems are configurations of the same workflow (Table~\ref{tab:systems}).
\begin{table}[t]
\caption{Evaluated systems (all use the same code path).}
\label{tab:systems}
\centering
\small
\begin{tabular}{p{2.6cm}p{5.2cm}}
\toprule
System & Configuration \\
\midrule
Closed-book & No retrieval. \\
Naive RAG & Top-5 whole paragraphs, one answer call. \\
Iterative RAG & Naive RAG plus Critic loop (no routing). \\
Adaptive-RAG (no CE) & Analyzer routing, type-dependent $k$, whole paragraphs, Critic loop. \\
\textbf{ContextRouter} & Adaptive-RAG plus Context Engineer with per-type budget. \\
CR w/o adaptive budget & Fixed $B{=}900$, $k{=}8$. \\
CR w/o critic loop & Single pass, no iteration. \\
\bottomrule
\end{tabular}
\end{table}

\textbf{Metrics.} Exact match (EM) and token F1 with the official HotpotQA answer normalisation; \emph{context recall}
(fraction of gold supporting paragraphs present in the final context); final \emph{context tokens}; mean LLM calls per
question; mean end-to-end LLM latency (sum of call latencies); and contract violations per question. We report 95\%
bootstrap confidence intervals (2000 resamples) and paired bootstrap intervals for differences between systems on
the same questions. \textbf{Budget sweep.} To isolate the effect of \emph{what} is placed in the window, we compare
the Context Engineer with hard truncation of rank-ordered paragraphs at identical budgets $B$ (single pass, no
routing).

\section{Results}
\textbf{Main comparison.} Table~\ref{tab:main} and Fig.~\ref{fig:main} report all systems on \nQ\ questions.
ContextRouter obtains F1 \fOnecr\% (95\% CI [\fLocr, \fHicr]) and EM \emScr\% with \ctxTokcr\ mean context tokens and
\callscr\ LLM calls per question. The adaptive pipeline without context engineering obtains F1 \fOnear\% with
\ctxTokar\ context tokens and \callsar\ calls. The paired F1 difference of ContextRouter over it is \dFar\ points
(95\% CI [\dFloar, \dFhiar]); mean context size changes by \dTokPctar\%, total tokens (prompt plus completion over all
calls) by \dTotPctar\% and LLM calls by \dCallsPctar\%. Against naive RAG the paired F1 difference is \dFnv\ points
(95\% CI [\dFlonv, \dFhinv]) and against iterative RAG \dFit\ points (95\% CI [\dFloit, \dFhiit]).
An interval that excludes zero indicates a reliable difference at this sample size; an interval that contains zero
indicates that the data do not distinguish the two systems.

\begin{table}[t]
\caption{Main results on HotpotQA (distractor) subset. EM/F1/recall in \%; Ctx tok.\ = mean final context tokens.}
\label{tab:main}
\centering
\scriptsize
\setlength{\tabcolsep}{3pt}
\begin{tabular}{l}
\hline
(results pending: run scripts/evaluate.py then scripts/make\_report.py)\\
\hline
\end{tabular}

\end{table}

\begin{figure}[t]
\centering
\figorpending{figs/fig_main.png}{\columnwidth}
\caption{F1 with 95\% CI (left) and mean final context tokens (right).}
\label{fig:main}
\end{figure}

\textbf{Evidence coverage.} Context recall is \recallcr\% for ContextRouter, \recallar\% for Adaptive-RAG without context
engineering and \recallnv\% for naive RAG, so the table shows directly whether the compressed context keeps the
supporting paragraphs. Mean Critic-triggered iterations are \itersar\ (Adaptive-RAG) and \iterscr\ (ContextRouter) per
question, and the router sent \directcr\% of questions to the no-retrieval route.

\textbf{Ablations.} Removing the adaptive budget gives F1 \fOnecrfb\% with \ctxTokcrfb\ tokens; removing the Critic loop
gives F1 \fOnecrnc\% with \callscrnc\ calls (full system: \fOnecr\% and \callscr).

\textbf{Budget sweep.} Table~\ref{tab:sweep} and Fig.~\ref{fig:sweep} compare context engineering with hard truncation at
equal budgets. At $B{=}\swLowB$ tokens the F1 values are \swcrLow\% (engineering) and \swrawLow\% (truncation); at
$B{=}\swMidB$ they are \swcrMid\% and \swrawMid\%. Context engineering has the higher F1 at \swWins\ of \swTotal\
budgets, with a mean difference of \swMeanGain\ F1 points.

\begin{table}[t]
\caption{Accuracy vs.\ context budget: hard truncation vs.\ Context Engineer (F1 in \%).}
\label{tab:sweep}
\centering
\small
\begin{tabular}{l}
\hline
(sweep pending)\\
\hline
\end{tabular}

\end{table}

\begin{figure}[t]
\centering
\figorpending{figs/fig_sweep.png}{0.8\columnwidth}
\caption{F1 as a function of the context budget $B$.}
\label{fig:sweep}
\end{figure}

\textbf{Interface reliability.} The mean number of contract violations per question is \violcr\ for ContextRouter
(\violar\ for Adaptive-RAG), so the typed contracts can be audited directly: a value near zero means the models returned schema-valid messages almost always, and a larger value quantifies the interface failures that the repair retry had to absorb.

\section{Discussion}
The experiment is designed so that each comparison answers one question. ContextRouter versus Adaptive-RAG isolates
the Context Engineer, because routing, retrieval and critic are shared; the sweep isolates \emph{what} is packed from
\emph{how much}, because both arms obey the same $B$. Context recall is the diagnostic for the outcomes: engineering can only help
answer quality if it retains the supporting paragraphs while discarding distractor sentences, and it can only
save tokens if the budget is below the size of the raw top-$k$ paragraphs. When the paired intervals contain zero, the
defensible claim is cost reduction at statistically indistinguishable quality, not an accuracy gain; when they
exclude zero, the sign states the direction. The ablations show which components (adaptive budget, critic loop)
contribute and at what cost in LLM calls.

\textbf{Limitations.} (1) One benchmark (HotpotQA) and a subsample of \nQ\ questions, so confidence intervals are wide
and conclusions are limited to this setting; (2) the Context Engineer is a hand-weighted extractive heuristic whose
weights were not learned; (3) a lexical retriever limits first-hop recall, which constrains every system equally;
(4) token counts for context are estimated (1.3 tokens per word), whereas prompt and completion tokens are API-reported;
(5) a single LLM family served by one provider, and API non-determinism despite temperature 0; (6) correctness is
measured by EM/F1 against short gold answers, which under-credits valid paraphrases.

\textbf{Future work.} Learning the sentence-scoring weights, a dense or hybrid retriever, additional benchmarks
(2WikiMultihopQA, MuSiQue) and a domain QA set built from course material, an LLM-judged faithfulness metric
\cite{es2024ragas}, and a learned budget policy.

\section{Conclusion}
We introduced ContextRouter, an adaptive RAG workflow that treats context construction as a budgeted, measurable agent
responsibility and coordinates its agents through typed contracts. The released evaluation compares it with
closed-book, naive, iterative and adaptive baselines, ablations, and a budget sweep against hard truncation, with
paired confidence intervals so that the effect of context engineering on quality, context size and cost can be read
directly from the data.

\section*{Reproducibility}
Source code, the sampled question list (fixed seed), evaluation scripts, per-question outputs and the cached LLM
responses are provided in the submission archive; the dataset is the public HotpotQA distractor validation split
\cite{yang2018hotpotqa}. Numbers in this manuscript are generated from \texttt{results/*.csv} by
\texttt{scripts/make\_report.py}.

\bibliographystyle{IEEEtran}
\bibliography{refs}
\end{document}
